\section{Định hướng kiểm thử và đánh giá}
\label{sec:kiem-chung}
\subsection{Tình huống kiểm thử}
Kiểm thử được tổ chức theo các nhóm thủ tục trong phạm vi, gồm câu hỏi đã
đủ thông tin, câu hỏi thiếu hoặc mơ hồ và hội thoại có bổ sung, sửa dữ kiện.
Cần có cả tình huống một miền và nhiều miền, khác thời điểm áp dụng, nhiệm
vụ chuyển tiếp, giấy tờ tồn đọng, nguồn thiếu hoặc mâu thuẫn và yêu cầu
ngoài phạm vi. Các biến thể này giúp kiểm tra hành vi tư vấn trong những
điều kiện khác nhau thay vì chỉ thử một câu hỏi có sẵn đầy đủ dữ kiện.

Mỗi ca cần có câu hỏi, dữ kiện ban đầu, thông tin được cung cấp ở các lượt
sau, nguồn và phiên bản đúng, kết luận mong đợi, điểm cần hỏi hoặc dừng.
Tình huống được dùng ở dạng được phép hoặc đã ẩn danh; các ví dụ giả định
được nhận diện riêng. Đáp án trọng yếu cần được xác nhận chuyên môn trước
khi dùng để chấm kết quả.

\subsection{Tiêu chí đánh giá}
Các tiêu chí được gắn với F01--F05 để kiểm tra trực tiếp yêu cầu sản phẩm:
\begin{itemize}[leftmargin=*,itemsep=2pt,topsep=3pt]
\item Hướng dẫn đúng và đủ đối với dữ kiện đã cung cấp, không kết luận vượt
 phạm vi căn cứ.
\item Câu hỏi bổ sung thu được thông tin cần thiết; nhận diện cả thông tin
 bị bỏ sót và những câu hỏi thừa.
\item Nguồn và phiên bản được chọn đúng; đoạn trích hỗ trợ nội dung kết luận.
\item Phản hồi sau khi sửa dữ kiện dùng thông tin mới, cập nhật đúng phần
 bị ảnh hưởng và giữ phần còn phù hợp.
\item Phân biệt được phần đáp ứng, không đáp ứng và chưa đủ căn cứ; dừng đúng
 khi nguồn chưa hỗ trợ kết luận.
\end{itemize}
Khả năng sử dụng và thời gian phản hồi cũng cần được ghi nhận trên cấu hình
chạy cụ thể. Các chỉ tiêu và kết quả đo sẽ được trình bày trong Chương~2.

\subsection{Nguyên tắc kiểm thử và đối chiếu}
Các cấu hình dự kiến gồm C0 là hỏi đáp truy xuất cơ bản, C1 bổ sung hỏi làm
rõ và C2 bổ sung đối chiếu điều kiện cùng cập nhật kết luận. So sánh C0--C1
cho phép xem tác động của hỏi bổ sung; so sánh C1--C2 xem tác động của nhóm
cơ chế đối chiếu và cập nhật, chưa tách riêng từng cơ chế.

Các cấu hình dùng cùng nguồn, dữ kiện đầu, phản hồi người dùng đã chuẩn bị,
mô hình nền và ngân sách xử lý. Tập phát triển và tập kiểm thử cần được
phân chia theo tình huống gốc; các cách diễn đạt lại của cùng một ca không
được xem như các trường hợp độc lập để đánh giá khả năng tổng quát. Cấu
hình chạy, bộ ca và số liệu cụ thể được trình bày ở Chương~2.

\subsection{Điều kiện áp dụng và giới hạn đánh giá}
Trước khi đánh giá cần chốt phiên bản kho nguồn, phạm vi từng nhánh và trạng
thái xác nhận đáp án. Một nhánh biết thông báo thiếu căn cứ chưa được coi
là hoàn thành nội dung tư vấn nếu nguồn cần thiết vẫn thiếu. Việc cập nhật
văn bản sau thời điểm chốt kho thuộc công tác duy trì hệ thống.

Kết quả chỉ có ý nghĩa trong phạm vi nguồn, đối tượng và tập tình huống đã
kiểm tra. Việc ứng dụng được vào hỗ trợ thực tế còn phụ thuộc quyền sử dụng
tài liệu, bảo vệ thông tin và xác nhận chuyên môn. Hiệu quả giảm thời gian
hoặc công việc cần được đo trong điều kiện sử dụng tương ứng.
